\section{Sharp fixed-level and modular refinements}
\label{sec:fixed-level-roadmap}

Proposition~\ref{prop:recursive-component-selection} identifies every
component on which contraction can remain trapped, and
Theorem~\ref{thm:simultaneous-marker-escape} handles its complement uniformly.
The retained fixed levels give exact classifications below the uniform field
threshold and resolve the first modular carriers.  Their proofs use stagewise polar
packages, marker and ramification divisors, and explicit adjacent-zero Hankel
kernels.  These calculations refine the general theorem; they are not its
induction steps.

The exceptional-orbit and split-free classifications at redundancies six and
seven are the fixed-level geometric contribution, stated in
Theorems~\ref{thm:r6} and~\ref{thm:r7}.  Their proofs separately supply the
covering-radius promotion; Proposition~\ref{prop:upper-radius} concerns the
higher fixed levels and the general large-characteristic range.

\begin{theorem}[Redundancy-six and redundancy-seven deep holes]
\label{thm:spine}%
\coverage{conditional_deduction}%
\lean{PRSRedundancySixSeven.redundancySixAllFieldSynthesis,redundancySevenAllFieldSynthesis,PRSPolarInduction.fifthPower_sigmaInversionOrbitCount}%
The following statements hold.
\begin{enumerate}[label=\textup{(\alph*)}]
\item For every prime power $q\geq7$, the projective deep-hole
      directions of $\PRS(q-5)$ are exactly the persistent tangent and
      conjugate-secant families, the five finite-field exception rows of
      Theorem~\ref{thm:r6}, and, when $q=2^m$ with odd $m\geq5$, one
      characteristic-two nucleus orbit.
\item For every prime power $q\geq7$, Theorem~\ref{thm:r7} gives the
      complete split-free redundancy-seven classification.  For $q\geq13$
      only the persistent families and, when $q=2^m$ with odd $m$, one
      central nucleus point remain.  The result is a deep-hole classification
      for every $q\geq7$.  The radius is six except at $q=8$, where it is
      seven; exact extraction
      leaves the nine-direction diagonal tangent orbit and the fixed central
      nucleus as the ten deep directions.
\end{enumerate}
\end{theorem}

At redundancies six and seven, the first two pointed packages expose the
stagewise mechanism in its smallest dimensions.  The persistent rank-two
carrier, its rank-one boundary, the modular nucleus, and the self-collision
divisor account for every polar line that can fail to escape.  The remaining
line meets a proper bad scheme of bounded degree, and the terminal cubic cover
supplies a split squarefree witness.  Appendix~\ref{sec:r67} proves the
exhaustive contained-component classification and closes the bounded fields
by finite certificates.  Higher fixed-level calculations are not needed for
the arbitrary-redundancy theorem and are left to the companion record.
